% ID: SEL-D0001-OB-E03
% ESTADO: validada. S002, 2026-11-07. Tres pasadas del mismo autor; ver informe.
Sean $S$ ver series y $P$ ver películas. Los datos son $P(S)=0.69$, $P(P)=0.35$ y $P(S^c\cap P^c)=0.18$.

\textbf{a) Series o películas.}\par\nopagebreak[4]

El suceso complementario de no ver ninguna es ver al menos una de las dos:
\[
P(S\cup P)=1-P(S^c\cap P^c)=1-0.18=\boxed{0.82}.
\]
La palabra ``o'' incluye a quienes ven ambas.

\textbf{b) Películas sabiendo que ve series.}\par\nopagebreak[4]

La fórmula de la unión permite recuperar la intersección, que no está dada directamente:
\[
P(S\cap P)=P(S)+P(P)-P(S\cup P)
=0.69+0.35-0.82=0.22.
\]
Condicionar a que ve series cambia el grupo de referencia a $S$:
\[
P(P\mid S)=\frac{P(P\cap S)}{P(S)}
=\boxed{\frac{22}{69}\approx0.318841}.
\]
La probabilidad es aproximadamente del $31.88\%$. No se necesita suponer independencia.

\textbf{c) Series y no películas.}\par\nopagebreak[4]

Dentro de quienes ven series distinguimos los que también ven películas y los que no:
\[
P(S\cap P^c)=P(S)-P(S\cap P)=0.69-0.22=\boxed{0.47}.
\]
\textbf{Comprobación.} Las probabilidades de ver ambas, sólo series, sólo películas o ninguna son, respectivamente,
\[
0.22,\qquad0.47,\qquad0.13,\qquad0.18.
\]
Son no negativas, suman 1 y reproducen los marginales $0.22+0.47=0.69$ y $0.22+0.13=0.35$.
